\documentclass[../main.tex]{subfiles}
\graphicspath{{\subfix{../tables_figures/}}}

\begin{document}
\section{Results}

\subsection{Health data description}

\begin{wrapfigure}{r}{0.5\textwidth}
  \centering
  \includegraphics[width=0.45\textwidth]{tables_figures/table1.png}
  \caption{Cases and controls by exposure status. PSPS event categories (mild/moderate/severe) are tertiles for each metric.}
  \label{table1}
\end{wrapfigure}

This zip code-level case-crossover study included 9,319,806 hospital admission or emergency department visits among individuals $\geq$ 20 years old for a respiratory, cardiovascular, or psychiatric endpoint between 2013 -- 2019 in California. Respiratory visits were most commonly observed (n = 4,923,047), followed by cardiovascular and psychiatric visits. 

PSPS events cover customers served by the three major energy providers in California -- Pacific Gas and Electric (PG\&E), San Diego Gas and Electric (SDG\&E), and Southern California Edison (SCE).\citep{cpuc_psps_reports} These three utilities serve approximately XX\% of California. Across these providers, we included XXXX PSPS events in this study. The median duration of a PSPS event was XX hours. There were XXXX zip code -days in our 7-year study period that experienced PSPS events, and each zip code experienced, on average, XXXX events (figure \ref{fig2}). Overall, XXXX of the YYYY zip codes in California experienced at least one PSPS event and were included in this study (figure \ref{fig1}).

Using these PSPS exposure data, we constructed two metrics -- an absolute and a relative metric. The absolute metric describes the number of customers impacted by a PSPS event and is a count. The relative metric describes YYY. We used tertiles to classify PSPS events as mild, moderate, and severe for each metric. Most PSPS events were classified as XXXX (Figure XXX). 
\begin{wrapfigure}{l}{0.5\textwidth}
    \includegraphics[width=\linewidth]{tables_figures/fig1.png}
    \caption{Zip codes included in study.}
    \label{fig1}
\end{wrapfigure}
We also examined wildfire smoke \PM both independent and concurrent with PSPS events. There were XXX zip code-days of wildfire smoke \PM from 2013 -- 2019 and YYY zip code-days of co-occurring wildfire smoke \PM and PSPS events during the same period (figure \ref{fig2}).

Using the HCAI data and our absolute PSPS metric, we identified XXX case-days and XX control-days across mild, moderate, and severe exposure for our absolute metric, and XXX case-days and XX control-days across mild, moderate, and severe exposure for our relative metric (table \ref{table1}). We identified the greatest number of encounters associated with the severe exposure level, with XXXX cardiovascular case-days during a severe outage and XXXX respiratory case-days (table \ref{table1}). Using our relative metric, we identified the greatest number of encounters associated with the XXX exposure level, with XXXX {cause} case-days during a severe outage and XXXX {cause} case-days (table \ref{table1}). 

\begin{figure}[htbp]
    \centering
    \includegraphics[width=\textwidth]{tables_figures/fig2.png}
    \caption{Zip codes included in study.}
    \label{fig2}
\end{figure}


\subsection{Respiratory morbidity}
 Independent association results first. Then joint. The odds of respiratory encounter were XX times higher in those exposed to co-occurring wildfire \PM and severe PSPS events (figure \ref{fig1}). NOTE: THIS FIGURE IS WRONG, IT IS A PLACEHOLDER!!!!

When we examined chronic obstructive pulmonary disease (COPD) alone, the effects attenuated and were null. This is likely due to small sample size. 

\begin{figure}[htbp]
    \centering
    % Remove extra white space with trim option
    % Adjust these values based on your specific image: {left} {bottom} {right} {top}
    \includegraphics[width=\textwidth,trim={60pt 40pt 60pt 40pt},clip]{tables_figures/fig3.png}
    \caption{Zip codes included in study.}
    \label{fig3}
\end{figure}

\subsection{Cardiovascular mobidity and Psychiatric disorders}
All findings for the association between PSPS events and psychiatric disorders were null. Additionally, results for the association between PSPS events that co-occur with wildfire \PM were also null. Though previous studies have found a relationship between wildfire \PM and psychiatric disorders such as dementia,\citep{elser2025wildfire} we did not have a large enough sample size to detect this relationship in isolation (\ref{fig1}). 

\end{document}